\subsection{Time Stepping Methods}
\label{SEC:Time:Methods}

GFR exclusively employs explicit Runge-Kutta (ERK) time-stepping methods for
advancing simulations in time.
%
There are currently 28 different ERK methods available in GFR; 6 of these
methods are recommended due to significantly greater performance compared to the
other available methods.
%
These 28 methods can be separated into two different types of ERK methods:
traditional ERK methods that use the Courant-Friedrichs-Lewy (CFL) number to
determine the optimal time step size that maintains stability, and embedded-pair
ERK (EPERK) methods that use a Proportional-Integral-Derivative (PID) controller
to automate the selection of an optimal time step size.
%
Further details on the ERK methods in GFR can be found in
\linecite{Spiegel_2025}.

%%%%%%%%%%%%%%%%%%%%%%%%%%%%%%%%%%%%%%%%%%%%%%%%%%%%%%%%%%%%%%%%%%%%%%%%%%%%%%%%

Each of the 28 ERK methods available in GFR are listed within one of the
following three sections where the method is categorized as
\quoterm{Recommended}, \quoterm{Other}, and \quoterm{\emph{NOT} Recommended}.
%
The CFL-based and EPERK methods are listed separately within each category for
easier identification.
%
Each method listed includes:
the \inp[nmlTimestep::RungeKuttaScheme]{Runge\_Kutta\_Scheme} input value
to activate it;
the number of Runge-Kutta stages, order, and description of the scheme; and
the speedup relative to the 3-stage/\nth{3}-order SSP ERK method (SSP3) which
is used as a baseline for performance.
%
%\begin{itemize}
%% \item[$\diamond$]
%  \item[]
%       [Input Value] : $i$-stage/$j$-order Description$^{\text{Ref}}$
%         (Speedup relative to SSP3)
%\end{itemize}
%%
%where
%[Input Value] refers to the value of 
%\inp[nmlTimestep::RungeKuttaScheme]{Runge\_Kutta\_Scheme}, and
%SSP3 refers to the 3-stage/\nth{3}-order SSP ERK method which is used as a
%baseline for performance.
%
It should be noted that the speedup referenced here is in terms of
time-to-solution, or the physical\slash real-world time taken to complete just
the solver part of a simulation.
%
Additionally, care was taken to ensure that simulations used to evaluate speedup
were identical aside from the ERK method used to advance the simulations in
time, \forexample all simulations were run on the same system with the same
parallel settings with the only difference in the input file being the specified
ERK method.
%
Speedup data is not available for every method.

%%%%%%%%%%%%%%%%%%%%%%%%%%%%%%%%%%%%%%%%%%%%%%%%%%%%%%%%%%%%%%%%%%%%%%%%%%%%%%%%

\subsubsection{Highly Recommended Runge-Kutta Methods}
\label{SEC:Time:Methods:Recommended}

The methods in \tref{TAB:RKMethods_Good} are highly recommended due to
significantly greater performance compared to the baseline SSP3 method.
%
Although these methods require the evaluation of $1.7\times$--$5\times$ as many
right-hand-side evaluations (\thatis number of stages), the maximum stable CFL
allowed by each of these methods compared to that of the baseline SSP3 method 
more than offsets this additional computational cost resulting in significant
reductions in time-to-solution.

\FINALFIX{\newpage}

%%%%%%%%%%%%%%%%%%%%%%%%%%%%%%%%%%%%%%%%%%%%%%%%%%%%%%%%%%%%%%%%%%%%%%%%%%%%%%%%

\begin{table}[h!]
\begin{center}
\begin{threeparttable}
%-------------------------------------------------------------------------------
\caption{Recommended Runge-Kutta Methods.}
\label{TAB:RKMethods_Good}
%-------------------------------------------------------------------------------
% Set S column format to right-align integers and then center in column.
\sisetup{
  table-alignment-mode = format,
  table-number-alignment = center,
  table-figures-decimal = 0,
  table-figures-integer = 3
}
\begin{tabular}{ S r @{ / } l l c }
  \toprule
  \multicolumn{1}{l}{\inp[nmlTimestep::RungeKuttaScheme]{Runge\_Kutta\_Scheme}}
        &    Stages &         Order & Description                & 
  Speedup\tnote{\ensuremath{\ast}} \\
  \midrule\\[-1.00em]
  \multicolumn{5}{l}{CFL-Based Methods} \\
    113 &  11-stage & \nth{3}-order & VOERK\cite{AlJahdali_2022} & 1.33x \\
    154 &  15-stage & \nth{4}-order & VOERK\cite{AlJahdali_2022} & 1.33x \\
      4 &   5-stage & \nth{4}-order & SSP ERK\cite{Spiteri_2002} & 1.22x \\
    165 &  16-stage & \nth{5}-order & VOERK\cite{AlJahdali_2022} & 1.18x \\
  \midrule\\[-1.00em]
  \multicolumn{5}{l}{Adaptive Time Stepping Methods} \\
    943 &   9-stage & \nth{4}(\nth{3})-order
                    & EPERK FSAL\cite{Ranocha_2022} & 1.22x \\
    532 &   5-stage & \nth{3}(\nth{2})-order
                    & EPERK FSAL\cite{Ranocha_2022} & 1.09x \\
  \bottomrule
\end{tabular}
%-------------------------------------------------------------------------------
\begin{tablenotes}
  \item[\ensuremath{\ast}]
  Relative to \inptext{Runge\_Kutta\_Scheme}=3,
  the 3-stage/\nth{3}-order SSP ERK method.
\end{tablenotes}
%-------------------------------------------------------------------------------
\end{threeparttable}
\end{center}
\end{table}

%%%%%%%%%%%%%%%%%%%%%%%%%%%%%%%%%%%%%%%%%%%%%%%%%%%%%%%%%%%%%%%%%%%%%%%%%%%%%%%%



%%%%%%%%%%%%%%%%%%%%%%%%%%%%%%%%%%%%%%%%%%%%%%%%%%%%%%%%%%%%%%%%%%%%%%%%%%%%%%%%

\subsubsection{Other Runge-Kutta Methods}
\label{SEC:Time:Methods:Other}

The methods in \tref{TAB:RKMethods_Other} are categorized as \quoterm{other}
because they are neither high- nor low-performing with similar performance to
the baseline SSP3 method.
%
This section contains the most popular ERK methods used in CFD, particularly the
Classic RK family of methods, SSP2, the 5-stage/\nth{4}-order Carpenter-Kennedy
low-storage RK method, and, probably the most popular of all ERK methods
especially with FR\slash DG, SSP3.

%%%%%%%%%%%%%%%%%%%%%%%%%%%%%%%%%%%%%%%%%%%%%%%%%%%%%%%%%%%%%%%%%%%%%%%%%%%%%%%%
\begin{table}[h!]
\begin{center}
\begin{threeparttable}
%-------------------------------------------------------------------------------
\caption{Other Runge-Kutta Methods.}
\label{TAB:RKMethods_Other}
%-------------------------------------------------------------------------------
% Set S column format to right-align integers and then center in column.
\sisetup{
  table-alignment-mode = format,
  table-number-alignment = center,
}
%  table-figures-decimal = 0,
%  table-figures-integer = 4
%}
\begin{tabular}{ S[table-format = +3] r @{ / } l l c }
  \toprule
  \multicolumn{1}{l}{\inp[nmlTimestep::RungeKuttaScheme]{Runge\_Kutta\_Scheme}}
        &    Stages &         Order & Description                & 
  Speedup\tnote{\ensuremath{\ast}} \\
  \midrule\\[-1.00em]
  \multicolumn{5}{l}{CFL-Based Methods} \\
     -2 &   2-stage & \nth{2}-order & Classic RK                     & n/a   \\
     -3 &   3-stage & \nth{3}-order & Classic RK                     & n/a   \\
     -4 &   4-stage & \nth{4}-order & Classic RK                     & n/a   \\
      2 &   2-stage & \nth{2}-order & SSP ERK\cite{Gottlieb_1998}    & n/a   \\
      3 &   3-stage & \nth{3}-order & SSP ERK\cite{Gottlieb_1998}    & 1.00x \\
      5 &   5-stage & \nth{4}-order & Carpenter-Kennedy              & 0.99x \\
     \multicolumn{3}{l}{}           & low storage RK\cite{Carpenter_1994} &  \\
     53 &   5-stage & \nth{3}-order & VOERK\cite{AlJahdali_2022}     & 1.08x \\
     64 &   6-stage & \nth{4}-order & VOERK\cite{AlJahdali_2022}     & 1.02x \\
     85 &   8-stage & \nth{5}-order & VOERK\cite{AlJahdali_2022}     & 1.05x \\
  \midrule\\[-1.00em]
  \multicolumn{5}{l}{Adaptive Time Stepping Methods} \\
   -943 &   9-stage & \nth{4}(\nth{3})-order
                    & EPERK non-FSAL\cite{Ranocha_2022} & n/a   \\
   -532 &   5-stage & \nth{3}(\nth{2})-order
                    & EPERK non-FSAL\cite{Ranocha_2022} & n/a   \\
  \bottomrule
\end{tabular}
%-------------------------------------------------------------------------------
\begin{tablenotes}
  \item[\ensuremath{\ast}]
  Relative to \inptext{Runge\_Kutta\_Scheme}=3,
  the 3-stage/\nth{3}-order SSP ERK method.
\end{tablenotes}
%-------------------------------------------------------------------------------
\end{threeparttable}
\end{center}
\end{table}

%%%%%%%%%%%%%%%%%%%%%%%%%%%%%%%%%%%%%%%%%%%%%%%%%%%%%%%%%%%%%%%%%%%%%%%%%%%%%%%%



%%%%%%%%%%%%%%%%%%%%%%%%%%%%%%%%%%%%%%%%%%%%%%%%%%%%%%%%%%%%%%%%%%%%%%%%%%%%%%%%

\subsubsection{\emph{NOT} Recommended Runge-Kutta Methods}
\label{SEC:Time:Methods:Bad}

The methods in \tref{TAB:RKMethods_Bad} are not recommended mostly due to
significantly worse performance compared to the baseline SSP3 method.
%
The one exception is the 1-stage/\nth{1}-order Classic RK method (\thatis
Forward Euler method), which is not recommended due to its low accuracy and
small CFL limit.
%
Although some of the AOERK methods achieve similar time-to-solution performance
compared to SSP3, they have all been listed here as not recommended because
every one of these methods is worse than their respective VOERK method
counterpart.

%%%%%%%%%%%%%%%%%%%%%%%%%%%%%%%%%%%%%%%%%%%%%%%%%%%%%%%%%%%%%%%%%%%%%%%%%%%%%%%%

\begin{table}[h!]
\begin{center}
\begin{threeparttable}
%-------------------------------------------------------------------------------
\caption{Not Recommended Runge-Kutta Methods.}
\label{TAB:RKMethods_Bad}
%-------------------------------------------------------------------------------
% Set S column format to right-align integers and then center in column.
\sisetup{
  table-alignment-mode = format,
  table-number-alignment = center,
}
%  table-figures-decimal = 0,
%  table-figures-integer = 5
%}
\begin{tabular}{ S[table-format = +4] r @{ / } l l c }
  \toprule
  \multicolumn{1}{l}{\inp[nmlTimestep::RungeKuttaScheme]{Runge\_Kutta\_Scheme}}
        &    Stages &         Order & Description                & 
  Speedup\tnote{\ensuremath{\ast}} \\
  \midrule\\[-1.00em]
  \multicolumn{5}{l}{CFL-Based Methods} \\
   \pm 1 &   1-stage & \nth{1}-order & Forward Euler method       & n/a \\
     -53 &   5-stage & \nth{3}-order & AOERK\cite{AlJahdali_2022} & n/a \\
    -113 &  11-stage & \nth{3}-order & AOERK\cite{AlJahdali_2022} & n/a \\
     -64 &   6-stage & \nth{4}-order & AOERK\cite{AlJahdali_2022} & n/a \\
    -154 &  15-stage & \nth{4}-order & AOERK\cite{AlJahdali_2022} & n/a \\
     -85 &   8-stage & \nth{5}-order & AOERK\cite{AlJahdali_2022} & n/a \\
    -165 &  16-stage & \nth{5}-order & AOERK\cite{AlJahdali_2022} & 0.98x \\
  \midrule\\[-1.00em]
  \multicolumn{5}{l}{Adaptive Time Stepping Methods} \\
    54 &   6-stage & \nth{5}(\nth{4})-order
                   & Dormand-Prince EPERK\cite{Dormand_1980} & 0.64x \\
   853 &  12-stage & \nth{8}(\nth{5})-order
                   & Dormand-Prince EPERK\cite{Hairer_1993}  & 0.61x \\
  1054 &  10-stage & \nth{5}(\nth{4})-order
                   & EPERK FSAL\cite{Ranocha_2022}           & 0.85x \\
 -1054 &  10-stage & \nth{5}(\nth{4})-order
                   & EPERK non-FSAL\cite{Ranocha_2022}       & n/a \\
  \bottomrule
\end{tabular}
%-------------------------------------------------------------------------------
\begin{tablenotes}
  \item[\ensuremath{\ast}]
  Relative to \inptext{Runge\_Kutta\_Scheme}=3,
  the 3-stage/\nth{3}-order SSP ERK method.
\end{tablenotes}
%-------------------------------------------------------------------------------
\end{threeparttable}
\end{center}
\end{table}

%%%%%%%%%%%%%%%%%%%%%%%%%%%%%%%%%%%%%%%%%%%%%%%%%%%%%%%%%%%%%%%%%%%%%%%%%%%%%%%%



%%%%%%%%%%%%%%%%%%%%%%%%%%%%%%%%%%%%%%%%%%%%%%%%%%%%%%%%%%%%%%%%%%%%%%%%%%%%%%%%

